% !TeX program = xelatex
% Компиляция: xelatex C46_classical_1000K.tex
% Исходник самодостаточен; данные таблиц включены в текст.
\documentclass[11pt,a4paper]{article}
\usepackage{fontspec}
\usepackage[russian]{babel}
\IfFontExistsTF{Cambria}{\setmainfont{Cambria}}{\setmainfont{DejaVu Serif}}
\IfFontExistsTF{Calibri}{\setsansfont{Calibri}}{\setsansfont{DejaVu Sans}}
\IfFontExistsTF{Consolas}{\setmonofont[Scale=MatchLowercase]{Consolas}}{\setmonofont[Scale=MatchLowercase]{DejaVu Sans Mono}}
\usepackage[a4paper,left=23mm,right=21mm,top=20mm,bottom=19mm,headheight=15pt,headsep=6mm,footskip=10mm]{geometry}
\usepackage{amsmath,amssymb,booktabs,array,xcolor,titlesec,fancyhdr,caption,hyperref}
\definecolor{ink}{HTML}{213542}
\definecolor{teal}{HTML}{006D70}
\definecolor{muted}{HTML}{526673}
\hypersetup{unicode=true,colorlinks=true,linkcolor=teal,urlcolor=teal,
 pdftitle={Классическая статистика колебаний в C46 для Ti–Al при 1000 К},
 pdfsubject={Равновесные составы фаз α и α₂ при фиксированном отношении осей}}
\urlstyle{same}
\captionsetup{font=small,labelfont={bf,color=teal},labelsep=period,singlelinecheck=false}
\titleformat{\section}{\sffamily\bfseries\Large\color{ink}}{\color{teal}\thesection}{.5em}{}
\titleformat{\subsection}{\sffamily\bfseries\large\color{ink}}{\thesubsection}{.5em}{}
\titlespacing*{\section}{0pt}{12pt plus 2pt}{5pt}
\titlespacing*{\subsection}{0pt}{9pt plus 2pt}{4pt}
\setlength{\parindent}{1.1em}
\setlength{\parskip}{2.4pt plus .6pt}
\setlength{\emergencystretch}{2em}
\brokenpenalty=10000 \widowpenalty=8000 \clubpenalty=8000
\renewcommand{\arraystretch}{1.16}
\raggedbottom
\pagestyle{fancy}\fancyhf{}
\fancyhead[L]{\small\sffamily\color{muted}Ti--Al\enspace/\enspace Классические колебания в C46}
\fancyhead[R]{\small\sffamily\color{muted}\(1000\) К}
\fancyfoot[R]{\small\sffamily\color{muted}\thepage}
\renewcommand{\headrulewidth}{.25pt}
\begin{document}
\thispagestyle{plain}
\noindent{\sffamily\small\color{teal}\bfseries ДОПОЛНЕНИЕ К НАУЧНОМУ ОТЧЁТУ\hfill 25 ноября 2026 г.\par}
\vspace{3mm}
\noindent{\sffamily\bfseries\fontsize{23}{26}\selectfont\color{ink}
Классические колебания в модели C46\par}
\vspace{1mm}
\noindent{\sffamily\large\color{muted}Границы \(\alpha/\alpha_2\) в Ti--Al при 1000~К и нулевом давлении\par}
\vspace{2mm}\noindent{\color{teal}\rule{\linewidth}{1pt}}\par\vspace{1mm}
\noindent Кластерное разложение C46 с классической статистикой колебаний даёт
\textbf{14,90 ат.\% Al в \(\alpha\)} и
\textbf{21,92 ат.\% Al в \(\alpha_2\)}.
На 5120 узлах этот результат получен как на 13 окнах оригинала, так и на
уточнённой сетке из 25 окон. Условные 95\%-е интервалы на густой сетке
составляют \(14{,}861\)--\(14{,}933\) и \(21{,}867\)--\(21{,}978\)
ат.\% Al. Классический и квантовый варианты основаны на
одних декорациях C46 и статических энергиях потенциала Qi и соавторов.
В обоих расчётах объём флуктуирует, а \(c/a=1{,}56374\) фиксировано.

\subsection*{Сдвиг обеих границ при замене статистики}
На исходных 13 окнах квантовая модель даёт \(14{,}943/21{,}931\) ат.\%
Al. Архивные значения C46, \(14{,}958/21{,}958\), входят в интервалы
нового квантового расчёта. Переход к классике сдвигает границу
\(\alpha\) к меньшим содержаниям Al. Для \(\alpha_2\) интервал парной разности
включает нуль. Поле сосуществования уменьшается на \(2{,}015\)~мэВ на
замену Ti на Al. В таблице сдвиг отсчитан от квантовой границы, найденной
на том же наборе траекторий.


\begin{center}\small
\begin{tabular}{@{}rr rr rr rr@{}}
\toprule
 && \multicolumn{2}{c}{Квантовые, ат.\%}
 & \multicolumn{2}{c}{Классические, ат.\%}
 & \multicolumn{2}{c}{Сдвиг, п.п.}\\
\cmidrule(lr){3-4}\cmidrule(lr){5-6}\cmidrule(l){7-8}
\(N\) & Окна & \(\alpha\) & \(\alpha_2\) & \(\alpha\) & \(\alpha_2\) & \(\alpha\) & \(\alpha_2\)\\
\midrule
768 & 13 & $14{,}982$ & $22{,}022$ & $14{,}913$ & $22{,}018$ & $-0{,}068$ & $-0{,}004$\\
5120 & 13 & $14{,}943$ & $21{,}931$ & $14{,}896$ & $21{,}922$ & $-0{,}047$ & $-0{,}009$\\
5120 & 25 & $14{,}933$ & $21{,}922$ & $14{,}896$ & $21{,}921$ & $-0{,}037$ & $-0{,}001$\\
\midrule
\multicolumn{8}{l}{95\%-е интервалы парного сдвига, п.п.}\\
768 & 13 & \multicolumn{3}{c}{\(\alpha\)\enspace $[-0{,}083;\,-0{,}052]$} & \multicolumn{3}{c}{\(\alpha_2\)\enspace $[-0{,}017;\,+0{,}010]$}\\
5120 & 13 & \multicolumn{3}{c}{\(\alpha\)\enspace $[-0{,}058;\,-0{,}036]$} & \multicolumn{3}{c}{\(\alpha_2\)\enspace $[-0{,}023;\,+0{,}005]$}\\
5120 & 25 & \multicolumn{3}{c}{\(\alpha\)\enspace $[-0{,}048;\,-0{,}027]$} & \multicolumn{3}{c}{\(\alpha_2\)\enspace $[-0{,}015;\,+0{,}012]$}\\
\bottomrule
\end{tabular}\end{center}


При переходе от 13 к 25 окнам абсолютные классические границы изменились
меньше чем на \(0{,}002\) п.п. На 25 окнах замена интерполяции PCHIP
линейными отрезками даёт парные сдвиги \(-0{,}036\) и \(-0{,}003\) п.п.
При уменьшении ячейки до 768 узлов сдвиг границы \(\alpha\) сохраняет
знак и возрастает по модулю. Для \(\alpha_2\) при обоих размерах
парный сдвиг сопоставим со статистической погрешностью.

\subsection*{Сравнение с прямым классическим MTP}
При тех же \(T,p,c/a\) расчёт с явными смещениями на 1400 атомах при
\(W<0{,}40\) дал \(16{,}083\) и \(22{,}159\) ат.\% Al. Относительно
классического C46 на авторской сетке 13 окон расхождение
«прямой MTP минус C46» составляет \(+1{,}187\) п.п. для \(\alpha\)
и \(+0{,}237\) п.п. для \(\alpha_2\). Для уточнённой сетки получены
\(+1{,}187\;[1{,}123;\,1{,}257]\) и \(+0{,}238\;[0{,}131;\,0{,}336]\)
п.п. Интервалы учитывают статистику двух независимых расчётов.
Поправка за статистику ядер слишком мала, чтобы объяснить расхождение
на Ti-богатой границе. В остаток входят усечение кластерного разложения,
квазигармоническое приближение, ограничения на смещения прямого MTP,
конечные размеры и способ извлечения составов.
\clearpage
\subsection*{Что измеряет парный расчёт Монте-Карло}

В объёмно-зависимом кластерном разложении C46 на каждой конфигурации
заполнения \(\sigma\) и удельном объёме \(v\) известны суммы по орбитам
\(S_\gamma(\sigma)\) для всей ячейки. Полная эффективная энергия
\(H_m(\sigma,v)=\sum_\gamma S_\gamma J_{\gamma,m}(v,T)\) включает
статическую часть MTP и колебательную свободную энергию для выбранной
статистики \(m\). На девяти узлах по объёму подгоняются все 20
коэффициентов, между узлами используется кубический сплайн.
В Монте-Карло меняются химические типы атомов и общий объём
при неизменной гексагональной форме. Внешнее давление равно нулю.

Чтобы измерить производную свободной энергии по составу \(x=N_{\rm Al}/N\),
в окне \(l\) к \(H_m\) добавляли известное смещение
\(-N\mu_0 x+Nk_{\rm B}T\kappa(x-x_{0,l})^2\).
Для обеих моделей использованы \(\kappa=100\) и
\(\mu_0=3{,}37431251\)~эВ на замену Ti на Al. По средним составам
определяли химическое поле и общую касательную,
\[
 \widehat{\Delta\mu}_{m,l}
 =\mu_0-2k_{\rm B}T\kappa(\langle x\rangle_{m,l}-x_{0,l}),\qquad
 \int_{x_\alpha}^{x_{\alpha_2}}
 [\widehat{\Delta\mu}_{m}(x)-\Delta\mu^*_m]\,dx=0 .
\]
Поля на концах интегрирования равны \(\Delta\mu_m^*\). Изотерму между
окнами интерполировали кусочно-кубической функцией с сохранением формы
(PCHIP). Положение двух контактов зависит от дискретизации изотермы,
поэтому основной контроль её разрешения выполнен удвоением плотности окон.

В 13 общих окнах центры \(x_0\) менялись от 0,13 до 0,25 с шагом
0,01. Сохраняли состав, объём и все суммы по орбитам. Это позволяет
точно вычислить \(\Delta H=H_{\rm cl}-H_{\rm q}\) на каждом кадре,
включая константу и точечную орбиту в текущем объёме.
Для данной ячейки и окна химическое смещение одинаково в обеих моделях.
Отношение весов классической и квантовой моделей пропорционально
\(\exp[-\Delta H/(k_{\rm B}T)]\), где в экспоненте стоит энергия
всей ячейки. Двусторонняя оценка Беннетта объединяет траектории
обеих моделей в одном окне. Благодаря общим кадрам оценка разности
границ сохраняет корреляцию между двумя рассчитанными изотермами.

Для \(N=5120\) в каждом из 13 окон одна серия начиналась в
упорядоченной конфигурации с классическими коэффициентами, другая~---
в случайной конфигурации с квантовыми. На серию приходилось 5000
проходов уравновешивания и 20\,000 производственных проходов, с
записью каждого десятого. В дополнительные 12 окон с центрами от
0,135 до 0,245 с шагом 0,01 запускали классические случайные старты
той же длины. В этих окнах квантовые средние получены обратным
перевзвешиванием с множителем \(\exp[+\Delta H/(k_{\rm B}T)]\).
В каждой бутстрэп-реализации для обеих моделей брали одни и те же
временные блоки из каждой траектории и заново находили обе общие
касательные. Интервал разности получали по распределению парных
разностей составов. Чувствительность к числу узлов, размеру ячейки
и выбору интерполяции оценивается отдельно.

\paragraph{Проверка перекрытия и начальных состояний.}
В 13 исходных окнах \(N=5120\) минимальная доля эффективных кадров
Киша после перевзвешивания составила \(0{,}904\); она измеряет
концентрацию весов до поправки на временную корреляцию. Число
временных блоков менялось от 19 до 200 на отдельную траекторию.
По квантовым случайным и классическим упорядоченным стартам раздельно
сдвиг \(x_\alpha\) равен \(-0{,}039\) и \(-0{,}057\)~п.п.,
а сдвиг \(x_{\alpha_2}\) равен \(+0{,}006\) и \(-0{,}012\)~п.п.
Абсолютный правый контакт классической модели по этим двум сериям
равен \(22{,}020\) и \(21{,}840\) ат.\% Al. Разность \(0{,}180\)~п.п.
имеет условный 95\%-й интервал \([0{,}049;\,0{,}294]\)~п.п. и составляет
2,87 суммарной стандартной ошибки. Для разностей парных поправок
между этими сериями 95\%-е интервалы включают нуль на обеих границах.
Увеличение длины блоков вдвое сохранило выводы по парным сдвигам.
Для 13 узлов обе границы разрешены во всех 768 совместных
бутстрэп-реализациях; при одновременной проверке 13 и 25 окон~---
во всех 512 реализациях. Разность полей для густой сетки составляет
\(-1{,}928\)~мэВ на замену.
\clearpage
\subsection*{Переход от квантовых фононов C46 к классическим}

Сохранены авторские свободные энергии 46 структур на девяти объёмах и
температурах от 0 до 2200~К с шагом 25~К. Полный комплект исходных
спектров для этих структур отсутствует в архиве. Для каждой устойчивой
пары «структура~--- объём» из полной энергии вычли статическую метку
и восстановили классическую колебательную энергию при \(T_0=1000\)~К.
Для одной гармонической моды частоты \(\omega\) при
\(z=\hbar\omega/(k_{\rm B}T)\) квантовая свободная энергия равна
\(k_{\rm B}T\ln[2\sinh(z/2)]\), классическая~--- \(k_{\rm B}T\ln z\).
Их разность начинается с \(k_{\rm B}T(z^2/24-z^4/2880)\).
После суммирования по модам и нормировки на атом поэтому получаем
\[
 \frac{F_{\rm q}^{\rm vib}(T)}{T/T_0}
 = F_{\rm cl}^{\rm vib}(T_0)
   -g k_{\rm B}T_0\ln(T/T_0)
   +\sum_{j=1}^{4} b_j(T_0/T)^{2j}
   +O\big((T_0/T)^{10}\big).
\]
Искомая энергия~--- свободный член. Вес \(g\) сохранённых мод на атом
определяли вместе с коэффициентами \(b_j\), чтобы учесть обработку
поступательных мод в авторской фононной сетке. При подгонке на
1000--2200~К и проверке с нижними границами 800, 1200 и 1400~К и
порядками 5, 4 и 3 максимальный разброс \(F_{\rm cl}^{\rm vib}\)
составил \(2{,}4\cdot10^{-12}\)~эВ/атом по всем 409 устойчивым парам.
Для 396 устойчивых спектров из независимой выборки ATAT ту же процедуру
проверили прямой суммой по частотам. Максимальная ошибка составила
\(7{,}0\cdot10^{-13}\)~эВ/атом. Восемь химических заполнений совпали
с авторскими. По всем девяти объёмам среднеквадратическое расхождение
между прямыми классическими энергиями ATAT и реконструкцией C46 для этих
структур составляет \(0{,}0136\)~мэВ/атом, при \(V_0\)~---
\(0{,}00104\)~мэВ/атом.

На каждом объёме классические метки подогнали в исходном базисе из
20 орбит с прежними весами, масками обучения и устойчивости.
Статические взаимодействия, отношение \(c/a=1{,}56374\) и обработка
объёмной зависимости сохранены. Поправка
\(\delta f=F_{\rm cl}^{\rm vib}-F_{\rm q}^{\rm vib}\) по всей выборке
лежит между \(-1{,}217\) и \(-0{,}491\)~мэВ/атом. При опорном объёме
её наилучшая линейная по составу часть по 46 структурам составляет
\(-0{,}7212-0{,}6661x_{\rm Al}\)~мэВ/атом, где \(x_{\rm Al}\)~---
атомная доля. Среднеквадратический остаток равен \(0{,}0466\), максимальный
по модулю~--- \(0{,}1251\)~мэВ/атом. При фиксированном объёме линейный
по составу член меняет наклон общей касательной и сохраняет составы.
В QHA все члены удержаны, поскольку их зависимость от объёма влияет на
минимум свободной энергии.

\begin{center}
\small
\begin{tabular}{@{}lrrrrrr@{}}
\toprule
 & \multicolumn{2}{c}{Квантовые, ат.\% Al}
 & \multicolumn{2}{c}{Классические, ат.\% Al}
 & \multicolumn{2}{c}{Сдвиг, п.п.}\\
\cmidrule(lr){2-3}\cmidrule(lr){4-5}\cmidrule(l){6-7}
Среднее поле & \(\alpha\) & \(\alpha_2\) & \(\alpha\) & \(\alpha_2\)
 & \(\alpha\) & \(\alpha_2\)\\
\midrule
HA  & 8,926 & 19,421 & 8,902 & 19,424 & \(-0{,}023\) & \(+0{,}003\)\\
QHA & 8,789 & 21,308 & 8,754 & 21,319 & \(-0{,}035\) & \(+0{,}011\)\\
\bottomrule
\end{tabular}
\end{center}
В приближении среднего поля (MFA) переход к классической статистике
расширяет двухфазный интервал QHA при \(p=0\) на \(0{,}046\) п.п.
HA рассчитана при фиксированном \(V_0/N=17{,}28833\)~\AA$^3$/атом.
\subsection*{Воспроизведение}
Расчётный каталог проекта \nolinkurl{phase_reproduction/classical_c46_1000/}
содержит скрипт \nolinkurl{fit_classical.py} для перевода авторских меток в
классический предел, полученные сетки \nolinkurl{quantum_grid.npz} и
\nolinkurl{classical_grid.npz}, поструктурные значения
\nolinkurl{structure_labels.csv} и проверку по прямым фононным спектрам
\nolinkurl{spectra_validation.json}. Границы среднего поля рассчитаны на этих
двух сетках с граничными условиями сплайна \emph{not-a-knot}.
Результаты Монте-Карло содержатся в \nolinkurl{paired5120.json},
\nolinkurl{paired768_final.json} и \nolinkurl{dense_grid.json};
\nolinkurl{direct_comparison.json} содержит остаток относительно прямого MTP.

Исходные взаимодействия MTP опубликованы
\href{https://doi.org/10.1103/PhysRevMaterials.7.103602}{Qi и соавт. (2023)}.
Ансамбль с ограниченной дисперсией состава для сплавов рассмотрен
\href{https://doi.org/10.1103/PhysRevB.86.134204}{Садигом и Эрхартом (2012)},
двусторонняя оценка свободных энергий~---
\href{https://doi.org/10.1016/0021-9991(76)90078-4}{Беннеттом (1976)}.
\end{document}
